\section{Praktikum}
\label{sec:praktikum-bab13}

\begin{praktikumbox}
\vspace{1.1em}
\textbf{Judul:} Menelusuri percabangan dan perulangan pada program
array.\\[0.3em]
\textbf{Sub-CPMK:} 3.4 \quad\textbullet\quad
\textbf{Alokasi waktu:} 100 menit \quad\textbullet\quad
\textbf{Bentuk:} individu
\end{praktikumbox}

\subsection{Tujuan Praktikum}

Setelah menyelesaikan praktikum ini, mahasiswa mampu menelusuri
percabangan dengan \asm{CMP} dan \asm{Jcc}, perulangan dengan
\asm{LOOP}, serta perulangan bersarang pada pengolahan array.

\subsection{Alat dan Bahan}

\begin{itemize}
    \item Komputer dengan emu8086 terpasang.
    \item DOSBox dengan TASM, TLINK, dan DEBUG.
    \item Berkas \file{loop-maksimum.asm} dan
          \file{bubble-sort.asm} dari \file{code/asm/}.
    \item Lembar pengamatan register dan \eng{flag}.
\end{itemize}

\subsection{Langkah Kerja}

\begin{enumerate}
    \item \textbf{Menguji pencarian maksimum.}
    \begin{enumerate}
        \item Rakit \file{loop-maksimum.asm} dengan TASM, lalu link
              dengan TLINK.
        \item Jalankan program dan catat keluarannya.
        \item Buka program pada emu8086 dan jalankan
              \eng{single step} sampai label \code{CARI}.
        \item Catat \reg{AX}, \reg{SI}, \reg{CX}, dan \flag{CF}
              setiap kali \lstasm|CMP AX, [SI]| selesai.
        \item Jelaskan kapan \lstasm|JAE LANJUT| diambil dan kapan
              tidak.
    \end{enumerate}

    \item \textbf{Mengubah data maksimum.}
    \begin{enumerate}
        \item Ubah array menjadi \code{25, 7, 120, 41, 93, 18, 60, 5}.
        \item Rakit ulang dan jalankan program.
        \item Pastikan nilai maksimum berubah menjadi 120.
    \end{enumerate}

    \item \textbf{Menguji pengurutan gelembung.}
    \begin{enumerate}
        \item Rakit dan jalankan \file{bubble-sort.asm}.
        \item Catat keluarannya.
        \item Telusuri satu lintasan luar dan catat nilai \reg{CX},
              \reg{SI}, dan \reg{DI}.
        \item Jelaskan peran \lstasm|JBE TIDAK_TUKAR| pada
              pembandingan.
    \end{enumerate}

    \item \textbf{Mengubah data pengurutan.}
    \begin{enumerate}
        \item Ubah array menjadi \code{9, 8, 7, 6, 5, 4, 3, 2}.
        \item Rakit ulang dan jalankan program.
        \item Pastikan hasilnya terurut menaik.
    \end{enumerate}

    \item \textbf{Membandingkan LOOP dan DEC/JNZ.}
    \begin{enumerate}
        \item Ganti \lstasm|LOOP BANDING| dengan pasangan
              \lstasm|DEC CX| dan \lstasm|JNZ BANDING|.
        \item Rakit ulang dan pastikan hasil pengurutan tetap sama.
        \item Jelaskan perbedaan penggunaan \reg{CX} pada kedua
              bentuk.
    \end{enumerate}

    \item \textbf{Menulis percabangan if--else.}
    \begin{enumerate}
        \item Tulis potongan program yang membandingkan \reg{AX} dan
              \reg{BX}.
        \item Bila \reg{AX} lebih besar, isi \code{hasil} dengan 1;
              bila tidak, isi dengan 0.
        \item Gunakan \asm{CMP}, \asm{JLE}, dan \asm{JMP}.
    \end{enumerate}
\end{enumerate}

\subsection{Hasil yang Dikumpulkan}

\begin{enumerate}
    \item Berkas \file{loop-maksimum.asm} dan
          \file{bubble-sort.asm} versi modifikasi.
    \item Tangkapan layar keluaran kedua program.
    \item Tabel pengamatan \reg{AX}, \reg{SI}, \reg{CX}, dan
          \eng{flag}.
    \item Penjelasan pemilihan \asm{JAE} dan \asm{JBE}.
    \item Potongan program \code{if--else}.
\end{enumerate}

\begin{asesmenbox}
\vspace{1.1em}
\textbf{Asesmen praktikum 11 (bobot 5\% dari komponen praktikum).}

\begin{tabular}{@{}p{6.2cm}cp{4.6cm}@{}}
    \toprule
    \textbf{Indikator} & \textbf{Bobot} & \textbf{Kriteria} \\
    \midrule
    Ketepatan penelusuran percabangan & 30\% &
    Keputusan \asm{JAE}/\asm{JBE} dijelaskan dengan benar \\
    \addlinespace
    Ketepatan penelusuran perulangan & 30\% &
    Nilai \reg{CX}, \reg{SI}, dan \reg{DI} tercatat benar \\
    \addlinespace
    Program modifikasi berjalan & 25\% &
    Keluaran berubah sesuai data baru \\
    \addlinespace
    Ketepatan potongan if--else & 15\% &
    Alur dan label disusun dengan benar \\
    \bottomrule
\end{tabular}
\end{asesmenbox}
